\section{Kiểm thử}
\label{sec:kiem-thu}

\paragraph{Môi trường.} WSL2 (Linux 6.18, 18 CPU), Go 1.27.3, loopback. Ngày đo: 12/11/2026.
\begin{center}
  \small
  \begin{tabularx}{\linewidth}{lY}
    \texttt{go vet ./...} & sạch \\
    \texttt{go test -race ./...} & sạch (protocol, transport, server) \\
    \texttt{BTL\_LONG=1 go test -run Long -v ./internal/transport/} & các phép đo ở \cref{tab:do-transport}
  \end{tabularx}
\end{center}

\subsection{Giả lập mạng}
\label{subsec:gia-lap}
Loopback không bao giờ mất gói và có RTT dưới 1 ms, nên không kiểm tra được gì về độ tin cậy hay điều khiển tắc
nghẽn. \fname{LossyConn} bọc socket UDP thật và làm xấu chiều gửi theo cấu hình:

\begin{table}[H]
  \centering
  \caption{Các tham số của bộ giả lập mạng \fname{LossyConn}}
  \label{tab:lossy}
  \begin{tabularx}{\linewidth}{lY}
    \toprule
    \textbf{Tham số} & \textbf{Ý nghĩa} \\
    \midrule
    \texttt{Loss} & Bỏ gói ngẫu nhiên với xác suất cho trước: mất gói \emph{không} do nghẽn (sóng yếu) \\
    \texttt{Delay}, \texttt{Jitter} & Trễ đường truyền, dao động ngẫu nhiên nhưng giữ FIFO như hàng đợi router:
    gói sau không bao giờ tới trước gói trước \\
    \texttt{Reorder} & Cố ý cho mỗi gói trễ độc lập, tức đảo thứ tự hàng loạt \\
    \texttt{Rate}, \texttt{QueueBytes} & Nút cổ chai băng thông với bộ đệm; đầy thì bỏ gói ở đuôi: mất gói
    \emph{do nghẽn}, đúng loại tín hiệu mà CUBIC được thiết kế để phản ứng \\
    \bottomrule
  \end{tabularx}
\end{table}

Ứng dụng client/server cũng nhận cờ \texttt{-loss} để demo trên máy thật.

\subsection{Test đơn vị và tích hợp của tầng giao thức và transport}
65 test (20 ở \fname{internal/protocol}, 45 ở \fname{internal/transport}) chạy trong khoảng 15 giây:

\begin{longtable}{L{0.24\linewidth} L{0.7\linewidth}}
  \caption{Các nhóm test của tầng giao thức và transport}\label{tab:unit-test} \\
  \toprule
  \textbf{Nhóm} & \textbf{Nội dung kiểm tra} \\
  \midrule
  \endfirsthead
  \toprule
  \textbf{Nhóm} & \textbf{Nội dung kiểm tra} \\
  \midrule
  \endhead
  Định dạng gói, frame & Ghi rồi đọc lại mọi loại frame và header ra đúng giá trị; nhiều frame trong một gói;
  dữ liệu cụt hoặc type lạ trả lỗi, không panic; CLOSE reason bị cắt đúng ranh giới ký tự UTF-8 \\
  Handshake, meta file & Khung \texttt{[len][JSON]} liên tiếp; từ chối độ dài quá lớn; meta file đọc được khi
  dữ liệu tới từng byte một \\
  \fname{SeqBuffer} & Sắp lại tin đến sai thứ tự (cả hoán vị ngẫu nhiên); bỏ tin trùng; bỏ qua chỗ trống vĩnh
  viễn để không treo \\
  ACK & Khoảng SACK đúng với các lỗ hổng; gói trùng; tối đa 8 khoảng; giới hạn lịch sử \\
  RTT, RTO & Mẫu đầu, công thức cập nhật, RTO không tiến về RTT khi đường truyền đều, ngưỡng mất theo thời gian \\
  Ghép stream & Đúng thứ tự, đảo thứ tự, trùng, chồng lấn, FIN tới trước dữ liệu; cắt frame $\le$ 1000 byte;
  ghi vào stream đã đóng trả lỗi \\
  Mã hóa & Hai bên ra cùng khóa; sửa bất kỳ byte nào thì giải mã thất bại; nonce không trùng giữa hai chiều;
  không thấy plaintext trên dây (\cref{sec:bao-mat}) \\
  CUBIC & Slow start gấp đôi mỗi RTT; giảm còn 70\% khi mất; hình dạng đường bậc 3 quanh $W_\text{max}$; vùng
  Reno-friendly; sàn và trần \\
  Truyền qua UDP thật & 5 MB qua loss 20\% mỗi chiều: đúng từng byte; 1 MB qua loss 10\% với CUBIC; chiều server
  $\to$ client \\
  Vòng đời kết nối & \texttt{Close} báo đầu kia ngay; đóng \fname{Mux} đóng mọi kết nối; bên kia chết thì phát
  hiện sau $\approx$ 10 s \\
  Phát hiện mất, gửi lại & RTO backoff nhân đôi khi mạng chết; mất theo số đếm không tăng backoff; gửi lại theo
  đúng thứ tự gửi; ACK cũ không ghi đè cửa sổ mới \\
  Bộ giả lập & Jitter giữ FIFO; nút cổ chai bỏ gói khi bộ đệm đầy \\
  \bottomrule
\end{longtable}

\subsection{Phép đo tầng transport}
\label{subsec:do-transport}

\begin{table}[H]
  \centering
  \caption{Truyền một stream dài qua mạng giả lập (đo ngày 12/11/2026)}
  \label{tab:do-transport}
  \small
  \begin{tabularx}{\linewidth}{YlYr}
    \toprule
    \textbf{Kịch bản} & \textbf{CC} & \textbf{Kết quả} & \textbf{Gói mất} \\
    \midrule
    5 MB, loss 20\% mỗi chiều & cố định & 0,11 s, đúng từng byte & 22,4\% \\
    1 MB, loss 10\% & CUBIC & 1,0 s, đúng từng byte & 18,1\% \\
    50 MB, loss 10\%, jitter 20 ms & cố định 64 KB & 31,0 s, \textbf{1,69 MB/s} & 10,0\% \\
    5 MB, loss 10\%, jitter 20 ms & CUBIC & 48,3 s, 0,11 MB/s (\cref{subsec:cubic-gioi-han}) & 13,3\% \\
    20 MB, loss 0\%, jitter 20 ms & CUBIC & 1,06 s, 19,7 MB/s & 0,2\% \\
    60 MB, nút cổ chai 2 MB/s, RTT 40 ms & CUBIC & 31,3 s, \textbf{2,01 MB/s} (hết băng thông) & 0,4\% \\
    20 MB, nút cổ chai 2 MB/s, RTT 40 ms & cố định 64 KB & 14,3 s, 1,47 MB/s & 0\% \\
    \bottomrule
  \end{tabularx}
\end{table}

``Gói mất'' là tỉ lệ gói bên gửi coi là mất và đã gửi lại, so với tổng số gói đã gửi; nó lớn hơn tỉ lệ loss đặt
cho chiều dữ liệu vì ACK ở chiều ngược lại cũng bị mất. Với 20 MB qua jitter 20 ms mà không mất gói, chỉ 0,2\%
gói bị coi là mất oan (trước khi sửa bộ giả lập: 89\%, thử thách~\ref{tt:jitter}). Độ trễ tin chat khi đang
truyền file: \cref{tab:hol}.

\subsection{Test hệ thống}
\phanB{Bảng 10 kịch bản chạy server và client thật: (1) 3 client chat chung + DM; (2) loss 10\%, 20\%, jitter;
(3) file 1 KB / 10 MB / 200 MB, kiểm tra SHA-256; (4) gửi file lớn trong khi chat, độ trễ chat; (5) Ctrl+C /
\texttt{kill -9} một client; (6) tắt server rồi bật lại; (7) sai mật khẩu / user lạ / đăng nhập trùng; (8) gói
rác vào cổng UDP; (9) 20 client cùng lúc: số goroutine, CPU, bộ nhớ; (10) Wireshark.}
